\IfFileExists{stacks-project.cls}{%
\documentclass{stacks-project}
}{%
\documentclass{amsart}
}

% For dealing with references we use the comment environment
\usepackage{verbatim}
\newenvironment{reference}{\comment}{\endcomment}
%\newenvironment{reference}{}{}
\newenvironment{slogan}{\comment}{\endcomment}
\newenvironment{history}{\comment}{\endcomment}

% For commutative diagrams we use Xy-pic
\usepackage[all]{xy}

% We use 2cell for 2-commutative diagrams.
\xyoption{2cell}
\UseAllTwocells

% We use multicol for the list of chapters between chapters
\usepackage{multicol}

% This is generally recommended for better output
\usepackage{lmodern}
\usepackage[T1]{fontenc}

% For cross-file-references
\usepackage{xr-hyper}

% Package for hypertext links:
\usepackage{hyperref}

% For any local file, say "hello.tex" you want to link to please
% use \externaldocument[hello-]{hello}
\externaldocument[introduction-]{introduction}
\externaldocument[conventions-]{conventions}
\externaldocument[sets-]{sets}
\externaldocument[categories-]{categories}
\externaldocument[topology-]{topology}
\externaldocument[sheaves-]{sheaves}
\externaldocument[sites-]{sites}
\externaldocument[stacks-]{stacks}
\externaldocument[fields-]{fields}
\externaldocument[algebra-]{algebra}
\externaldocument[brauer-]{brauer}
\externaldocument[homology-]{homology}
\externaldocument[derived-]{derived}
\externaldocument[simplicial-]{simplicial}
\externaldocument[more-algebra-]{more-algebra}
\externaldocument[smoothing-]{smoothing}
\externaldocument[modules-]{modules}
\externaldocument[sites-modules-]{sites-modules}
\externaldocument[injectives-]{injectives}
\externaldocument[cohomology-]{cohomology}
\externaldocument[sites-cohomology-]{sites-cohomology}
\externaldocument[dga-]{dga}
\externaldocument[dpa-]{dpa}
\externaldocument[sdga-]{sdga}
\externaldocument[hypercovering-]{hypercovering}
\externaldocument[schemes-]{schemes}
\externaldocument[constructions-]{constructions}
\externaldocument[properties-]{properties}
\externaldocument[morphisms-]{morphisms}
\externaldocument[coherent-]{coherent}
\externaldocument[divisors-]{divisors}
\externaldocument[limits-]{limits}
\externaldocument[varieties-]{varieties}
\externaldocument[topologies-]{topologies}
\externaldocument[descent-]{descent}
\externaldocument[perfect-]{perfect}
\externaldocument[more-morphisms-]{more-morphisms}
\externaldocument[flat-]{flat}
\externaldocument[groupoids-]{groupoids}
\externaldocument[more-groupoids-]{more-groupoids}
\externaldocument[etale-]{etale}
\externaldocument[chow-]{chow}
\externaldocument[intersection-]{intersection}
\externaldocument[pic-]{pic}
\externaldocument[weil-]{weil}
\externaldocument[adequate-]{adequate}
\externaldocument[dualizing-]{dualizing}
\externaldocument[duality-]{duality}
\externaldocument[discriminant-]{discriminant}
\externaldocument[derham-]{derham}
\externaldocument[local-cohomology-]{local-cohomology}
\externaldocument[algebraization-]{algebraization}
\externaldocument[curves-]{curves}
\externaldocument[resolve-]{resolve}
\externaldocument[models-]{models}
\externaldocument[functors-]{functors}
\externaldocument[equiv-]{equiv}
\externaldocument[pione-]{pione}
\externaldocument[etale-cohomology-]{etale-cohomology}
\externaldocument[proetale-]{proetale}
\externaldocument[relative-cycles-]{relative-cycles}
\externaldocument[more-etale-]{more-etale}
\externaldocument[trace-]{trace}
\externaldocument[crystalline-]{crystalline}
\externaldocument[spaces-]{spaces}
\externaldocument[spaces-properties-]{spaces-properties}
\externaldocument[spaces-morphisms-]{spaces-morphisms}
\externaldocument[decent-spaces-]{decent-spaces}
\externaldocument[spaces-cohomology-]{spaces-cohomology}
\externaldocument[spaces-limits-]{spaces-limits}
\externaldocument[spaces-divisors-]{spaces-divisors}
\externaldocument[spaces-over-fields-]{spaces-over-fields}
\externaldocument[spaces-topologies-]{spaces-topologies}
\externaldocument[spaces-descent-]{spaces-descent}
\externaldocument[spaces-perfect-]{spaces-perfect}
\externaldocument[spaces-more-morphisms-]{spaces-more-morphisms}
\externaldocument[spaces-flat-]{spaces-flat}
\externaldocument[spaces-groupoids-]{spaces-groupoids}
\externaldocument[spaces-more-groupoids-]{spaces-more-groupoids}
\externaldocument[bootstrap-]{bootstrap}
\externaldocument[spaces-pushouts-]{spaces-pushouts}
\externaldocument[spaces-chow-]{spaces-chow}
\externaldocument[groupoids-quotients-]{groupoids-quotients}
\externaldocument[spaces-more-cohomology-]{spaces-more-cohomology}
\externaldocument[spaces-simplicial-]{spaces-simplicial}
\externaldocument[spaces-duality-]{spaces-duality}
\externaldocument[formal-spaces-]{formal-spaces}
\externaldocument[restricted-]{restricted}
\externaldocument[spaces-resolve-]{spaces-resolve}
\externaldocument[formal-defos-]{formal-defos}
\externaldocument[defos-]{defos}
\externaldocument[cotangent-]{cotangent}
\externaldocument[examples-defos-]{examples-defos}
\externaldocument[algebraic-]{algebraic}
\externaldocument[examples-stacks-]{examples-stacks}
\externaldocument[stacks-sheaves-]{stacks-sheaves}
\externaldocument[criteria-]{criteria}
\externaldocument[artin-]{artin}
\externaldocument[quot-]{quot}
\externaldocument[stacks-properties-]{stacks-properties}
\externaldocument[stacks-morphisms-]{stacks-morphisms}
\externaldocument[stacks-limits-]{stacks-limits}
\externaldocument[stacks-cohomology-]{stacks-cohomology}
\externaldocument[stacks-perfect-]{stacks-perfect}
\externaldocument[stacks-introduction-]{stacks-introduction}
\externaldocument[stacks-more-morphisms-]{stacks-more-morphisms}
\externaldocument[stacks-geometry-]{stacks-geometry}
\externaldocument[moduli-]{moduli}
\externaldocument[moduli-curves-]{moduli-curves}
\externaldocument[examples-]{examples}
\externaldocument[exercises-]{exercises}
\externaldocument[guide-]{guide}
\externaldocument[desirables-]{desirables}
\externaldocument[coding-]{coding}
\externaldocument[obsolete-]{obsolete}
\externaldocument[fdl-]{fdl}
\externaldocument[index-]{index}

% Theorem environments.
%
\theoremstyle{plain}
\newtheorem{theorem}[subsection]{Theorem}
\newtheorem{proposition}[subsection]{Proposition}
\newtheorem{lemma}[subsection]{Lemma}

\theoremstyle{definition}
\newtheorem{definition}[subsection]{Definition}
\newtheorem{example}[subsection]{Example}
\newtheorem{exercise}[subsection]{Exercise}
\newtheorem{situation}[subsection]{Situation}

\theoremstyle{remark}
\newtheorem{remark}[subsection]{Remark}
\newtheorem{remarks}[subsection]{Remarks}

\numberwithin{equation}{subsection}

% Macros
%
\def\lim{\mathop{\mathrm{lim}}\nolimits}
\def\colim{\mathop{\mathrm{colim}}\nolimits}
\def\Spec{\mathop{\mathrm{Spec}}}
\def\Hom{\mathop{\mathrm{Hom}}\nolimits}
\def\Ext{\mathop{\mathrm{Ext}}\nolimits}
\def\SheafHom{\mathop{\mathcal{H}\!\mathit{om}}\nolimits}
\def\SheafExt{\mathop{\mathcal{E}\!\mathit{xt}}\nolimits}
\def\Sch{\mathit{Sch}}
\def\Mor{\mathop{\mathrm{Mor}}\nolimits}
\def\Ob{\mathop{\mathrm{Ob}}\nolimits}
\def\Sh{\mathop{\mathit{Sh}}\nolimits}
\def\NL{\mathop{N\!L}\nolimits}
\def\CH{\mathop{\mathrm{CH}}\nolimits}
\def\proetale{{pro\text{-}\acute{e}tale}}
\def\etale{{\acute{e}tale}}
\def\QCoh{\mathit{QCoh}}
\def\Ker{\mathop{\mathrm{Ker}}}
\def\Im{\mathop{\mathrm{Im}}}
\def\Coker{\mathop{\mathrm{Coker}}}
\def\Coim{\mathop{\mathrm{Coim}}}

% Boxtimes
%
\DeclareMathSymbol{\boxtimes}{\mathbin}{AMSa}{"02}

%
% Macros for moduli stacks/spaces
%
\def\QCohstack{\mathcal{QC}\!\mathit{oh}}
\def\Cohstack{\mathcal{C}\!\mathit{oh}}
\def\Spacesstack{\mathcal{S}\!\mathit{paces}}
\def\Quotfunctor{\mathrm{Quot}}
\def\Hilbfunctor{\mathrm{Hilb}}
\def\Curvesstack{\mathcal{C}\!\mathit{urves}}
\def\Polarizedstack{\mathcal{P}\!\mathit{olarized}}
\def\Complexesstack{\mathcal{C}\!\mathit{omplexes}}
% \Pic is the operator that assigns to X its picard group, usage \Pic(X)
% \Picardstack_{X/B} denotes the Picard stack of X over B
% \Picardfunctor_{X/B} denotes the Picard functor of X over B
\def\Pic{\mathop{\mathrm{Pic}}\nolimits}
\def\Picardstack{\mathcal{P}\!\mathit{ic}}
\def\Picardfunctor{\mathrm{Pic}}
\def\Deformationcategory{\mathcal{D}\!\mathit{ef}}
\usepackage{cleveref}
\crefname{originallemma}{original lemma}{original lemmas}
\Crefname{originallemma}{Original lemma}{Original lemmas}
\makeatletter
\newlabel{categories-lemma-properties-2-cat-cats}{{003F}{source}{}{sourcepointer-1}{}}
\newlabel{categories-lemma-properties-2-cat-cats@cref}{{[originallemma][0][]003F}{[1][0][]source}{}{}{}}
\makeatother
\crefname{originalreference}{original reference}{original references}
\Crefname{originalreference}{Original reference}{Original references}
\makeatletter
\newlabel{example-godement}{{0G5M}{source}{}{sourcepointer-2}{}}
\newlabel{example-godement@cref}{{[originalreference][0][]0G5M}{[1][0][]source}{}{}{}}
\makeatother
\crefname{originalreference}{original reference}{original references}
\Crefname{originalreference}{Original reference}{Original references}
\makeatletter
\newlabel{lemma-characterize-simplicial-object}{{016B}{source}{}{sourcepointer-3}{}}
\newlabel{lemma-characterize-simplicial-object@cref}{{[originalreference][0][]016B}{[1][0][]source}{}{}{}}
\makeatother
\title{The Godement construction: d and s satisfying the two identities make X a simplicial endofunctor with augmentation}
\author{The Stacks Project Authors}
\date{stacks-project:simplicial.tex:local-snapshot-20261012; GNU FDL 1.2 or later}
\begin{document}
\maketitle
\noindent\textit{Editorial note.} All delivered mathematical text is the author's own, copied from the pinned Stacks Project chapter file sources/stacks-simplicial/pinned-original-source.tex; no statement or proof step was rewritten, reordered or removed.\par\medskip
\noindent\textit{Editorial note.} The counted proof is the author's complete argument: the checks of (1)(b) and (1)(d) invoke 'the same kind of calculus as in case (1)(a)' (source lines 6505 and 6522) with the resulting simplified expressions stated in-span - the referenced computation is the fully displayed template above them and the outcome is deterministic; the two assumptions of the lemma are used exactly where the simplicial identities demand them ((1)(c) and (1)(e) at j = i).\par\medskip
\noindent\textit{Editorial note.} The label lemma-godement names the counted statement, and the proof environment immediately following it in the source is delivered verbatim as its proof.\par\medskip
\noindent\textit{Editorial note.} The counted statement is carried by the author's own lemma environment in the chapter source.\par\medskip
\noindent\textit{Allowed standard prerequisites (editorial).}\begin{itemize}
\item The Godement setup: comonad-like data d, s on Y. Same chapter, \textbackslash{}label\{example-godement\}. Lines 6442-6449: the endofunctor Y with natural transformations d : Y -> id and s : id -> Y.
\item Characterization of simplicial objects by the five face/degeneracy relations. Same chapter, \textbackslash{}label\{lemma-characterize-simplicial-object\}. Lines 6455-6458: relations (1)(a)-(e) to be verified.
\item Interchange of star and circle compositions. Categories chapter, \textbackslash{}label\{categories-lemma-properties-2-cat-cats\}. Lines 6466-6470: (a' circ a) star (b' circ b) = (a' star b') circ (a star b), the engine of every computation.
\end{itemize}\medskip
\ifdefined\hypertarget\hypertarget{sourcepointer-1}{}\fi
\noindent\textit{Source reference.} Original lemma 003F: sources/stacks-chow/tags.txt, source lines 132--132. Prerequisite pointer invoked at proof line(s) 6473; it names a standard fact of a sibling chapter that the counted argument uses, and every step of the counted proof built on it is retained verbatim in the delivered body.\par\medskip
\ifdefined\hypertarget\hypertarget{sourcepointer-2}{}\fi
\noindent\textit{Source reference.} Original reference 0G5M: sources/stacks-simplicial/pinned-original-source.tex, source lines 6410--6410. Prerequisite pointer invoked at proof line(s) 6444; it names a standard fact of the same chapter that the counted argument uses, and every step of the counted proof built on it is retained verbatim in the delivered body.\par\medskip
\ifdefined\hypertarget\hypertarget{sourcepointer-3}{}\fi
\noindent\textit{Source reference.} Original reference 016B: sources/stacks-simplicial/pinned-original-source.tex, source lines 435--435. Prerequisite pointer invoked at proof line(s) 6457; it names a standard fact of the same chapter that the counted argument uses, and every step of the counted proof built on it is retained verbatim in the delivered body.\par\medskip
\begin{lemma}
\label{lemma-godement}
In Example \ref{example-godement} if
$$
1_Y = (d \star 1_Y) \circ s = (1_Y \star d) \circ s
\quad\text{and}\quad
(s \star 1) \circ s = (1 \star s) \circ s
$$
then $X = (X_n, d^n_j, s^n_j)$ is a simplicial object in the category
of endofunctors of $\mathcal{C}$ and $d : X_0 = Y \to \text{id}_\mathcal{C}$
defines an augmentation.
\end{lemma}

\begin{proof}
To see that we obtain a simplicial object we have to check that the
relations (1)(a) -- (e) of Lemma \ref{lemma-characterize-simplicial-object}
are satisfied. We will use the short hand notation
$$
1_a = 1_{X_{a - 1}} = 1_Y \star \ldots \star 1_Y \quad (a\text{ factors})
$$
for $a \geq 0$. With this notation we have
$$
d^n_j = 1_j \star d \star 1_{n - j}
\quad\text{and}\quad
s^n_j = 1_j \star s \star 1_{n - j}
$$
We are repeatedly going to use the rule that for transformations
of funtors $a, a', b, b'$ we have
$(a' \circ a) \star (b' \circ b) = (a' \star b') \circ (a \star b)$
provided that the $\star$ and $\circ$ compositions in this formula
make sense, see
Categories, Lemma \ref{categories-lemma-properties-2-cat-cats}.

\medskip\noindent
Condition (1)(a) always holds (no conditions needed on $d$ and $s$).
Namely, let $0 \leq i < j \leq n + 1$. We have to show that
$d^n_i \circ d^{n + 1}_j = d^n_{j - 1} \circ d^{n + 1}_i$, i.e.,
$$
(1_i \star d \star 1_{n - i}) \circ
(1_j \star d \star 1_{n + 1 - j}) =
(1_{j - 1} \star d \star 1_{n + 1 - j}) \circ
(1_i \star d \star 1_{n + 1 - i})
$$
We can rewrite the left hand side as
\begin{align*}
& (1_i \star d \star 1_{j - i - 1} \star 1_{n + 1 - j}) \circ
(1_i \star 1_1 \star 1_{j - i - 1} \star d \star 1_{n + 1 - j}) \\
& =
1_i \star
\left((d \star 1_{j - i - 1}) \circ
(1_1 \star 1_{j - i - 1} \star d)\right) \star 1_{n + 1 - j} \\
& =
1_i \star d \star 1_{j - i - 1} \star d \star 1_{n + 1 - j} 
\end{align*}
The second equality is true because $d \circ 1_1 = d$ and
$1_{j - i} \circ (1_{j - i - 1} \star d) = 1_{j - i - 1} \star d$.
A similar computation gives the same result for the right hand side.

\medskip\noindent
We check condition (1)(b). Let $0 \leq i < j \leq n - 1$. We have to show
that $d^n_i \circ s^{n - 1}_j = s^{n - 2}_{j - 1} \circ d^{n - 1}_i$, i.e.,
$$
(1_i \star d \star 1_{n - i}) \circ
(1_j \star s \star 1_{n - 1 - j}) =
(1_{j - 1} \star s \star 1_{n - 1 - j}) \circ
(1_i \star d \star 1_{n - 1 - i})
$$
By the same kind of calculus as in case (1)(a) both sides simplify
to $1_i \star d \star 1_{j - i - 1} \star s \star 1_{n - j - 1}$.

\medskip\noindent
We check condition (1)(c). Let $0 \leq j \leq n - 1$. We have to show
$\text{id} = d^n_j \circ s^{n - 1}_j = d^n_{j + 1} \circ s^{n - 1}_j$, i.e.,
$$
1_n =
(1_j \star d \star 1_{n - j}) \circ
(1_j \star s \star 1_{n - 1 - j}) =
(1_{j + 1} \star d \star 1_{n - j - 1}) \circ
(1_j \star s \star 1_{n - 1 - j})
$$
This is easily seen to be implied by the first assumption of the lemma.

\medskip\noindent
We check condition (1)(d). Let $0 < j + 1 < i \leq n$. We have to show
$d^n_i \circ s^{n - 1}_j = s^{n - 2}_j \circ d^{n - 1}_{i - 1}$, i.e.,
$$
(1_i \star d \star 1_{n - i}) \circ
(1_j \star s \star 1_{n - 1 - j}) =
(1_j \star s \star 1_{n - 2 - j}) \circ
(1_{i - 1} \star d \star 1_{n - i})
$$
By the same kind of calculus as in case (1)(a) both sides simplify
to $1_j \star s \star 1_{i - j - 2} \star d \star 1_{n - i}$.

\medskip\noindent
We check condition (1)(e). Let $0 \leq i \leq j \leq n - 1$. We
have to show that
$s^n_i \circ s^{n - 1}_j = s^n_{j + 1} \circ s^{n - 1}_i$, i.e.,
$$
(1_i \star s \star 1_{n - i}) \circ
(1_j \star s \star 1_{n - 1 - j}) =
(1_{j + 1} \star s \star 1_{n - 1 - j}) \circ
(1_i \star s \star 1_{n - 1 - i})
$$
By the same kind of calculus as in case (1)(a) this reduces to
$$
(s \star 1_{j - i + 1}) \circ (1_{j - i} \star s) =
(1_{j - i + 1} \star s) \circ (s \star 1_{j - i})
$$
If $j = i$ this is exactly one of the two assumptions of the lemma.
For $j > i$ left and
right hand side both reduce to the equality $s \star 1_{j - i - 1} \star s$
by calculations similar to those we did in case (1)(a).

\medskip\noindent
Finally, in order to show that $d$ defines an augmentation we have
to show that $d \circ (1_1 \star d) = d \circ (d \star 1_1)$ which is
true because both sides are equal to $d \star d$.
\end{proof}

\end{document}
